% !PRINTED 445
\weberpage{445}
%<<CONT
in die Gleichungen ein, durch die der unge\"anderte Punkt bestimmt wird [\S.\,116, (3)], so ergiebt sich:

\[
\begin{array}{l}
(\beta+1)\,x_1 + \gamma\,\varepsilon^3 x_2 + \alpha\,\varepsilon^2 x_3 = 0,\\
\gamma\,\varepsilon^{-3} x_1 + (\alpha+1)\,x_2 + \beta\,\varepsilon^{-1} x_3 = 0,\\
\alpha\,\varepsilon^{-2} x_1 + \beta\,\varepsilon\, x_2 + (\gamma+1)\,x_3 = 0.
\end{array}
\]
Diese drei Gleichungen reduciren sich alle auf die eine, die man z. B. aus der ersten durch Multiplication mit $\beta$ und Anwendung von \S.\,115, (8) ableitet:

\weq{5}{\alpha\,\gamma\,x_1 + \beta\,\gamma\,\varepsilon^3 x_2 + \alpha\,\beta\,\varepsilon^2 x_3 = 0,}
und die eine gerade Linie darstellt. Diese Linie ist also eine \gesperrt{Hauptaxe}.

Solcher Hauptaxen erhalten wir 21, da wir 21 Substitutionen 2\textsuperscript{ten} Grades haben.

Man kann die Functionen, die, gleich Null gesetzt, die 21 Hauptaxen darstellen, aus (5) leicht bilden, wenn man die Function

\weq{6}{A = \alpha\,\gamma\,x_1 + \beta\,\gamma\,\varepsilon^3 x_2 + \alpha\,\beta\,\varepsilon^2 x_3}
durch die Substitutionen $\tau^r$, $\tau^r\chi$, $\tau^r\chi^2$ transformirt ($r = 0, 1, 2, 3, 4, 5, 6$). Man erh\"alt so die 21 Functionen:

\weq{7}{\begin{array}{l}
A_{1,r} = \alpha\,\gamma\,\varepsilon^r x_1 + \gamma\,\beta\,\varepsilon^{2r+3} x_2 + \alpha\,\beta\,\varepsilon^{4r+2} x_3\\
A_{2,r} = \beta\,\gamma\,\varepsilon^{r+3} x_1 + \alpha\,\beta\,\varepsilon^{2r+2} x_2 + \alpha\,\gamma\,\varepsilon^{4r} x_3\\
A_{3,r} = \alpha\,\beta\,\varepsilon^{r+2} x_1 + \alpha\,\gamma\,\varepsilon^{2r} x_2 + \beta\,\gamma\,\varepsilon^{4r+3} x_3,
\end{array}}
aus denen leicht zu ersehen ist, dass die 21 Hauptaxen wirklich alle von einander verschieden sind.

Wir fassen dies so zusammen:

\noindent 2.\ \gesperrt{Es geh\"oren 21 verschiedene Hauptaxen zu der Gruppe $G_{168}$, von denen jede in jede andere durch Substitutionen der Gruppe transformirt werden kann.}

Da es nur 21 Hauptaxen giebt, so muss jede von ihnen durch acht Substitutionen unge\"andert bleiben, und diese acht Substitutionen bilden eine Gruppe. Um diese Gruppe zu ermitteln, wenden wir auf den in (6) dargestellten Ausdruck $A$ die Substitution $\omega$ [\S.\,115, (5)] an, und dadurch geht er \"uber in

\weq{8}{\begin{array}{l}
\gamma\,(\alpha^2 + \beta^2\varepsilon^3 + \alpha\,\beta\,\varepsilon^2)\,x_1 + \beta\,(\alpha\,\gamma + \gamma^2\varepsilon^3 + \alpha^2\varepsilon^2)\,x_2\\
{}+ \alpha\,(\gamma^2 + \beta\,\gamma\,\varepsilon^3 + \beta^2\varepsilon^2)\,x_3.
\end{array}}
